% !TEX program = xelatex
\documentclass[12pt]{article}
\usepackage[paperwidth=612bp,paperheight=792bp,margin=0bp,noheadfoot]{geometry}
\usepackage{fontspec}
\usepackage{xeCJK}
\usepackage{tikz}
\pagestyle{empty}
\setlength{\parindent}{0pt}
\setlength{\parskip}{0pt}
\IfFontExistsTF{TeX Gyre Termes}{\setmainfont{TeX Gyre Termes}}{\setmainfont{Times New Roman}}
\IfFontExistsTF{Noto Serif CJK SC}{\setCJKmainfont{Noto Serif CJK SC}}{\setCJKmainfont{Songti SC}}
\newcommand{\QuestionTitle}{人机实验1 H-C5a：高祖本纪}
\newcommand{\DrawQuestionTitle}{%
\node[anchor=base west,inner sep=0pt,outer sep=0pt] at ([xshift=72bp,yshift=-34bp]current page.north west) {{\fontsize{14bp}{14bp}\selectfont\bfseries \QuestionTitle}};%
\node[anchor=base west,inner sep=0pt,outer sep=0pt] at ([xshift=72bp,yshift=-62bp]current page.north west) {{\fontsize{12bp}{12bp}\selectfont 用时：\underline{\hspace{80bp}}}};%
}
\begin{document}
% Auto-generated from 刘邦_chapter_with_numbers.tex
\thispagestyle{empty}
\begin{tikzpicture}[remember picture,overlay]
\DrawQuestionTitle
\node[anchor=base west,inner sep=0pt,outer sep=0pt] at ([xshift=72bp,yshift=-84bp]current page.north west) {{\fontsize{12bp}{12bp}\selectfont 【8.16.1】 秦二世三年，楚怀王看到项梁的军队被打垮了，心里恐惧，迁离盱台，建都彭}};
\node[anchor=base west,inner sep=0pt,outer sep=0pt] at ([xshift=72bp,yshift=-100bp]current page.north west) {{\fontsize{12bp}{12bp}\selectfont 城，合并吕臣、项羽的军队，亲自统率。}};
\node[anchor=base west,inner sep=0pt,outer sep=0pt] at ([xshift=72bp,yshift=-116bp]current page.north west) {{\fontsize{12bp}{12bp}\selectfont 【8.16.2】 以沛公任砀郡长，封为武安侯，统领砀郡的军队。}};
\node[anchor=base west,inner sep=0pt,outer sep=0pt] at ([xshift=72bp,yshift=-132bp]current page.north west) {{\fontsize{12bp}{12bp}\selectfont 【8.16.3】 封项羽为长安侯，号为鲁公。}};
\node[anchor=base west,inner sep=0pt,outer sep=0pt] at ([xshift=72bp,yshift=-148bp]current page.north west) {{\fontsize{12bp}{12bp}\selectfont 【8.16.4】 吕臣任司徒，他的父亲吕青作令尹。}};
\node[anchor=base west,inner sep=0pt,outer sep=0pt] at ([xshift=72bp,yshift=-164bp]current page.north west) {{\fontsize{12bp}{12bp}\selectfont 【8.17.1】 赵多次请求救援，楚怀王就以宋义为上将军，项羽为次将，范增为未将，北上}};
\node[anchor=base west,inner sep=0pt,outer sep=0pt] at ([xshift=72bp,yshift=-180bp]current page.north west) {{\fontsize{12bp}{12bp}\selectfont 救赵。}};
\node[anchor=base west,inner sep=0pt,outer sep=0pt] at ([xshift=72bp,yshift=-196bp]current page.north west) {{\fontsize{12bp}{12bp}\selectfont 【8.17.2】 命令沛公西出略地，打入关中。}};
\node[anchor=base west,inner sep=0pt,outer sep=0pt] at ([xshift=72bp,yshift=-212bp]current page.north west) {{\fontsize{12bp}{12bp}\selectfont 【8.17.3】 同将领们约定：先攻入关中的，就封在关中做王。}};
\node[anchor=base west,inner sep=0pt,outer sep=0pt] at ([xshift=72bp,yshift=-228bp]current page.north west) {{\fontsize{12bp}{12bp}\selectfont 【8.18.1】 这时候，秦军强盛，常常乘胜追击，众将领没有认为先入关的是有利的。}};
\node[anchor=base west,inner sep=0pt,outer sep=0pt] at ([xshift=72bp,yshift=-244bp]current page.north west) {{\fontsize{12bp}{12bp}\selectfont 【8.18.2】 唯独项羽痛恨秦打垮了项梁的军队，心中愤激，愿和沛公西进入关。}};
\node[anchor=base west,inner sep=0pt,outer sep=0pt] at ([xshift=72bp,yshift=-260bp]current page.north west) {{\fontsize{12bp}{12bp}\selectfont 【8.18.3】 怀王的老将都说：“项羽为人轻捷而凶猛，狡诈而残忍。项羽曾经攻打襄城，}};
\node[anchor=base west,inner sep=0pt,outer sep=0pt] at ([xshift=72bp,yshift=-276bp]current page.north west) {{\fontsize{12bp}{12bp}\selectfont 襄城没有留下一个活人，全都活埋了。所经过的地方，无不残杀毁灭。况且楚军多次进兵}};
\node[anchor=base west,inner sep=0pt,outer sep=0pt] at ([xshift=72bp,yshift=-292bp]current page.north west) {{\fontsize{12bp}{12bp}\selectfont 攻取，（没有获胜，）以前陈王、项梁都失败了。不如另派宽厚长者，以正义为号召，向}};
\node[anchor=base west,inner sep=0pt,outer sep=0pt] at ([xshift=72bp,yshift=-308bp]current page.north west) {{\fontsize{12bp}{12bp}\selectfont 西进发，把道理向秦父老兄弟讲清楚。秦父老兄弟苦于他们君主的统治很久了，现在如果}};
\node[anchor=base west,inner sep=0pt,outer sep=0pt] at ([xshift=72bp,yshift=-324bp]current page.north west) {{\fontsize{12bp}{12bp}\selectfont 真能得到宽厚长者去关中，不加欺凌暴虐，应该能够拿下关中。而今项羽剽悍，不可派遣}};
\node[anchor=base west,inner sep=0pt,outer sep=0pt] at ([xshift=72bp,yshift=-340bp]current page.north west) {{\fontsize{12bp}{12bp}\selectfont 。只有沛公向来是宽大长者，可以派遣。”}};
\node[anchor=base west,inner sep=0pt,outer sep=0pt] at ([xshift=72bp,yshift=-356bp]current page.north west) {{\fontsize{12bp}{12bp}\selectfont 【8.18.4】 终于没有答应项羽，而派遣沛公西进攻取秦地。}};
\node[anchor=base west,inner sep=0pt,outer sep=0pt] at ([xshift=72bp,yshift=-372bp]current page.north west) {{\fontsize{12bp}{12bp}\selectfont 【8.18.5】 收集陈王、项梁的散兵，路经砀，到达成阳，与杠里的秦军对垒，打败了秦军}};
\node[anchor=base west,inner sep=0pt,outer sep=0pt] at ([xshift=72bp,yshift=-388bp]current page.north west) {{\fontsize{12bp}{12bp}\selectfont 的两支部队。}};
\node[anchor=base west,inner sep=0pt,outer sep=0pt] at ([xshift=72bp,yshift=-404bp]current page.north west) {{\fontsize{12bp}{12bp}\selectfont 【8.22.8】 沛公派人与秦朝官吏巡行县城乡间，告谕百姓。}};
\node[anchor=base west,inner sep=0pt,outer sep=0pt] at ([xshift=72bp,yshift=-420bp]current page.north west) {{\fontsize{12bp}{12bp}\selectfont 【8.22.9】 秦地的百姓大为高兴，争先恐后地拿出牛羊酒食款待士兵。}};
\node[anchor=base west,inner sep=0pt,outer sep=0pt] at ([xshift=72bp,yshift=-436bp]current page.north west) {{\fontsize{12bp}{12bp}\selectfont 【8.22.10】 沛公又谦让不肯接受，说：“仓库的谷子很多，不缺乏，不愿破费百姓。”}};
\node[anchor=base west,inner sep=0pt,outer sep=0pt] at ([xshift=72bp,yshift=-452bp]current page.north west) {{\fontsize{12bp}{12bp}\selectfont 百姓更加高兴，唯恐沛公不做秦王。}};
\node[anchor=base west,inner sep=0pt,outer sep=0pt] at ([xshift=72bp,yshift=-468bp]current page.north west) {{\fontsize{12bp}{12bp}\selectfont 【8.23.1】 有人劝沛公说：“秦地比天下富足十倍，地势好。如今听说章邯投降了项羽，}};
\node[anchor=base west,inner sep=0pt,outer sep=0pt] at ([xshift=72bp,yshift=-484bp]current page.north west) {{\fontsize{12bp}{12bp}\selectfont 项羽就给了雍王的封号，称王于关中。现在即将来到关中就国，你沛公恐怕不能占有这个}};
\node[anchor=base west,inner sep=0pt,outer sep=0pt] at ([xshift=72bp,yshift=-500bp]current page.north west) {{\fontsize{12bp}{12bp}\selectfont 地方了。应赶快派兵把守函谷关，不让诸侯军进来，逐渐征集关中兵，以加强实力，抵抗}};
\node[anchor=base west,inner sep=0pt,outer sep=0pt] at ([xshift=72bp,yshift=-516bp]current page.north west) {{\fontsize{12bp}{12bp}\selectfont 诸侯兵。”}};
\node[anchor=base west,inner sep=0pt,outer sep=0pt] at ([xshift=72bp,yshift=-532bp]current page.north west) {{\fontsize{12bp}{12bp}\selectfont 【8.23.2】 沛公赞成他的计策，照着做了。}};
\node[anchor=base west,inner sep=0pt,outer sep=0pt] at ([xshift=72bp,yshift=-548bp]current page.north west) {{\fontsize{12bp}{12bp}\selectfont 【8.23.3】 汉元年十一月间，项羽果然率领诸侯军西进，想要入关，而关门闭着。}};
\node[anchor=base west,inner sep=0pt,outer sep=0pt] at ([xshift=72bp,yshift=-564bp]current page.north west) {{\fontsize{12bp}{12bp}\selectfont 【8.23.4】 听说沛公已经平定关中，大怒，派黥布等攻破了函谷关。}};
\node[anchor=base west,inner sep=0pt,outer sep=0pt] at ([xshift=72bp,yshift=-580bp]current page.north west) {{\fontsize{12bp}{12bp}\selectfont 【8.23.5】 十二月间，就到了戏水。}};
\node[anchor=base west,inner sep=0pt,outer sep=0pt] at ([xshift=72bp,yshift=-596bp]current page.north west) {{\fontsize{12bp}{12bp}\selectfont 【8.23.6】 沛公左司马曹无伤听说项王发怒，}};
\node[anchor=base west,inner sep=0pt,outer sep=0pt] at ([xshift=72bp,yshift=-612bp]current page.north west) {{\fontsize{12bp}{12bp}\selectfont 【8.23.7】 要攻打沛公，派人告诉项羽说：“沛公想要称王关中，令子婴为相，珍宝被他}};
\node[anchor=base west,inner sep=0pt,outer sep=0pt] at ([xshift=72bp,yshift=-628bp]current page.north west) {{\fontsize{12bp}{12bp}\selectfont 全部占有了。”}};
\node[anchor=base west,inner sep=0pt,outer sep=0pt] at ([xshift=72bp,yshift=-644bp]current page.north west) {{\fontsize{12bp}{12bp}\selectfont 【8.23.8】 打算以此求得封赏。}};
\node[anchor=base west,inner sep=0pt,outer sep=0pt] at ([xshift=72bp,yshift=-660bp]current page.north west) {{\fontsize{12bp}{12bp}\selectfont 【8.23.9】 亚父劝项羽进攻沛公。}};
\end{tikzpicture}
\null
\end{document}

